\documentclass{article}
\usepackage[utf8]{inputenc}
\title{Python}
\begin{document}
\maketitle

\section{Python}
Python é uma linguagem de programação de alto nível, criada por **Guido van Rossum**
e lançada em 1991. É conhecida pela sua sintaxe _simples e legível_.


\subsection{Por que aprender Python?}
Python é atualmente uma das linguagens mais populares do mundo. É utilizada em
áreas como desenvolvimento web, inteligência artificial e automação de tarefas.


\subsubsection{Principais Vantagens}
\begin{itemize}
  \item **Sintaxe simples** — fácil de aprender para iniciantes

  \item **Multiplataforma** — corre em Windows, Linux e macOS

  \item **Comunidade ativa** — milhares de bibliotecas disponíveis

  \item **Versátil** — usado em web, ciência de dados, IA e scripting

\end{itemize}
\subsubsection{Instalação}
Para instalar Python, visita o site oficial:


[Download Python](https://www.python.org/downloads "Python.org")


![Python Logo "The Python Logo"](https://www.python.org/static/community_logos/python-logo.png)


\subsection{Tipos de Dados Básicos}
Python suporta vários tipos de dados nativos. Os mais comuns são:


\begin{enumerate}
  \item **int** — números inteiros, como `42` ou `-7`

  \item **float** — números decimais, como `3.14`

  \item **str** — cadeias de texto, como `"Olá, Mundo!"`

  \item **bool** — valores lógicos: `True` ou `False`

  \item **list** — coleções ordenadas e mutáveis

\end{enumerate}
\subsection{Estruturas de Controlo}
As estruturas de controlo em Python são _intuitivas_ e dispensam chavetas.
Utilizam indentação para delimitar blocos de código.


\subsubsection{Condicionais e Ciclos}
\begin{itemize}
  \item Condicional `if`

  \item Executa um bloco se a condição for verdadeira

  \item Ciclo `for`

  \item Itera sobre qualquer objeto iterável

  \item Ciclo `while`

  \item Repete enquanto a condição for verdadeira

\end{itemize}
\subsection{Bibliotecas Populares}
A força do Python está no seu ecossistema. Algumas bibliotecas essenciais:


\begin{enumerate}
  \item **NumPy** — computação numérica e arrays multidimensionais

  \item **Pandas** — análise e manipulação de dados

  \item **Matplotlib** — visualização de dados em gráficos

  \item **Flask** — desenvolvimento de aplicações web _minimalistas_

  \item **TensorFlow** — aprendizagem automática e redes neuronais

\end{enumerate}
\subsection{Recursos para Aprender}
Existem muitos recursos online gratuitos para aprender Python. Abaixo alguns
dos mais recomendados pela comunidade.


\begin{itemize}
  \item [Documentação Oficial](https://docs.python.org/3 "Docs Python 3")

  \item [Real Python](https://realpython.com "Real Python Tutorials")

  \item [W3Schools Python](https://www.w3schools.com/python "W3Schools")

\end{itemize}
\subsubsection{Conclusão}
Python é uma escolha **excelente** para qualquer programador, seja iniciante
ou experiente. A sua _versatilidade_ e a enorme comunidade tornam-no uma
ferramenta indispensável no mundo tecnológico atual.

\end{document}
